% Terjemahan Bahasa Indonesia dari chapter1.tex baris 763--925.
% Sumber: Methods of Algebra, Volume 2, commit
% 9a5803ff77dd3257484cb177f851a73770a59dd3 (CC BY 4.0).
% stable-unit-id: o014.aljabr2.chapter1.filtered-inductive-limits

% segment-id: o014.li2.chapter1.filtered-inductive-limits.h001
\section{Limit Induktif Terfilter}\label{sec:filtered-indlim}
\hypertarget{o014.aljabr2.chapter1.filtered-inductive-limits}{ }
\enlargethispage{2\baselineskip}

% segment-id: o014.li2.chapter1.filtered-inductive-limits.p001
Salah satu fakta umum dalam analisis matematika ialah bahwa, jika suatu
barisan konvergen, limitnya dapat dihitung dari subbarisan mana pun. Dalam
teori kategori, gagasan ini diwujudkan oleh konsep kofinalitas. Kita mulai
dari keterhubungan.

% segment-id: o014.li2.chapter1.filtered-inductive-limits.d001
\begin{definisi}\index{kategori!terhubung (connected)}
	Misalkan $I$ suatu kategori. Jika $\Obj(I) \neq \emptyset$ dan, untuk setiap
	$i, i' \in \Obj(I)$, selalu terdapat $n \geq 1$ beserta morfisme-morfisme
	% segment-id: o014.li2.chapter1.filtered-inductive-limits.q001
	\[ i = i_0 \leftarrow i_1 \to i_2 \leftarrow i_3 \to i_4 \leftarrow \cdots \to i_{2n} = i' , \]
	maka $I$ disebut terhubung.
\end{definisi}

% segment-id: o014.li2.chapter1.filtered-inductive-limits.p002
Kita perkenalkan sebuah konsep bantu yang akan digunakan berulang kali dalam
beberapa bagian berikutnya.

% segment-id: o014.li2.chapter1.filtered-inductive-limits.d002
\begin{definisi}[Kategori koma, lihat {\cite[Definisi 2.4.7]{Li1}}]\label{def:comma-category}
	\index{kategori koma (comma category)}
	\index[sym1]{iHHi@$(i/H)$, $(H/i)$}
	\index[sym1]{Pii@$\Pi_{i/}$, $\Pi_{/i}$}
	Diberikan kategori $I$, $J$, dan funktor $H: J \to I$. Misalkan
	$i \in \Obj(I)$.
	% segment-id: o014.li2.chapter1.filtered-inductive-limits.l001
	\begin{itemize}
		% segment-id: o014.li2.chapter1.filtered-inductive-limits.i001
		\item Menurut definisi, objek kategori koma $(i/H)$ (atau $(H/i)$)
		adalah data $(j, i \to Hj)$ (atau $(j, Hj \to i)$), dengan
		$j \in \Obj(J)$ dan $i \to Hj$ (atau $Hj \to i$) merupakan morfisme
		di $I$. Morfisme dari data $(j, i \to Hj)$ ke
		$(j', i \to Hj')$ (atau dari $(j, Hj \to i)$ ke
		$(j', Hj' \to i)$) adalah morfisme $f: j \to j'$ di $J$ yang membuat
		diagram berikut komutatif:
		% segment-id: o014.li2.chapter1.filtered-inductive-limits.g001
		\[\begin{tikzcd}[row sep=small, column sep=small]
			& i \arrow[ld] \arrow[rd] & \\
			Hj \arrow[rr, "{Hf}"'] & & Hj'
		\end{tikzcd} \quad \text{atau} \quad \begin{tikzcd}[row sep=small, column sep=small]
			& i & \\
			Hj \arrow[ru] \arrow[rr, "{Hf}"'] & & Hj' \arrow[lu] .
		\end{tikzcd}\]
		Komposisi dan morfisme identitas didefinisikan sebagaimana biasa.
		% segment-id: o014.li2.chapter1.filtered-inductive-limits.i002
		\item Kedua kategori di atas masing-masing dilengkapi dengan funktor
		proyeksi $\Pi_{i/}: (i/H) \to J$ dan $\Pi_{/i}: (H/i) \to J$, yang
		memetakan objek $(j, i \to Hj)$ (atau $(j, Hj \to i)$) ke $j$ dan
		memetakan morfisme $f$ ke $f$.
		% segment-id: o014.li2.chapter1.filtered-inductive-limits.i003
		\item Setiap morfisme $i \to i'$ di $I$ secara alami menginduksi funktor
		$(i'/H) \to (i/H)$ dan $(H/i) \to (H/i')$, sehingga diagram berikut
		komutatif:
		% segment-id: o014.li2.chapter1.filtered-inductive-limits.g002
		\[\begin{tikzcd}[column sep=small, row sep=small]
			(i'/H) \arrow[rr] \arrow[rd, "{\Pi_{i'/}}"'] & & (i/H) \arrow[ld, "{\Pi_{i/}}"] \\
			& J &
		\end{tikzcd} \quad \begin{tikzcd}[column sep=small, row sep=small]
			(H/i) \arrow[rr] \arrow[rd, "{\Pi_{/i}}"'] & & (H/i') \arrow[ld, "{\Pi_{/i'}}"] . \\
			& J &
		\end{tikzcd}\]
	\end{itemize}
\end{definisi}

% segment-id: o014.li2.chapter1.filtered-inductive-limits.e001
\begin{contoh}\label{eg:comma-vs-cone}
	Tinjau funktor $\alpha: I \to \mathcal{C}$ sebagai objek di
	$\mathcal{C}^I$, dan tinjau pula funktor diagonal
	$\Delta: \mathcal{C} \to \mathcal{C}^I$ yang memetakan setiap objek ke
	funktor konstan yang bersesuaian. Menurut Definisi
	\ref{def:comma-category}, $(\alpha/\Delta)$ tepat merupakan kategori
	``kokonus'' yang diperkenalkan dalam Konvensi \ref{con:limit-cone}, sedangkan
	funktor $\Pi_{\alpha/}$ memetakan setiap kokonus ke puncaknya.
\end{contoh}

% segment-id: o014.li2.chapter1.filtered-inductive-limits.d003
\begin{definisi}\label{def:cofinal} \index{kofinal (cofinal)}
	Untuk funktor $H: J \to I$, jika $(i/H)$ terhubung untuk setiap
	$i \in \Obj(I)$, maka $H$ disebut \emph{kofinal}; jika $H$ merupakan
	funktor inklusi dari suatu subkategori, subkategori $J$ disebut kofinal.
\end{definisi}

% segment-id: o014.li2.chapter1.filtered-inductive-limits.p003
Mudah diverifikasi bahwa jika $I \to J$ dan $J \to K$ keduanya kofinal, maka
funktor komposit $I \to K$ juga kofinal. Pemeriksaannya diserahkan kepada
pembaca sebagai latihan.

% segment-id: o014.li2.chapter1.filtered-inductive-limits.p004
Hasil berikut menunjukkan bahwa limit dapat dihitung dengan membatasi diagram
sepanjang funktor kofinal. Ingat bahwa $H$ menginduksi
$H^{\opp}: J^{\opp} \to I^{\opp}$.

% segment-id: o014.li2.chapter1.filtered-inductive-limits.pr001
\begin{proposisi}\label{prop:cofinal-lim}
	Jika $H: J \to I$ kofinal, maka:
	% segment-id: o014.li2.chapter1.filtered-inductive-limits.l002
	\begin{itemize}
		% segment-id: o014.li2.chapter1.filtered-inductive-limits.i004
		\item untuk funktor $\alpha: I \to \mathcal{C}$ yang diberikan,
		$\varinjlim \alpha$ ada jika dan hanya jika $\varinjlim \alpha H$ ada;
		dalam hal ini terdapat isomorfisme kanonik
		$\varinjlim \alpha H \simeq \varinjlim \alpha$;
		% segment-id: o014.li2.chapter1.filtered-inductive-limits.i005
		\item secara dual, untuk funktor
		$\beta: I^{\opp} \to \mathcal{C}$ yang diberikan,
		$\varprojlim \beta$ ada jika dan hanya jika
		$\varprojlim \beta H^{\opp}$ ada; dalam hal ini terdapat isomorfisme
		kanonik $\varprojlim \beta H^{\opp} \simeq \varprojlim \beta$.
	\end{itemize}
\end{proposisi}
% segment-id: o014.li2.chapter1.filtered-inductive-limits.pf001
\begin{bukti}
	Kita hanya membahas $\varinjlim$. Dalam bahasa
	\S\ref{sec:limit-functor}, cukup dibuktikan bahwa funktor
	$(\alpha/\Delta) \to (\alpha H / \Delta)$ merupakan ekuivalensi; funktor
	ini memetakan kokonus beralas $\alpha$, yakni $(L, (f_i)_i)$, ke kokonus
	beralas $\alpha H$, yakni $(L, (f_{Hj})_j)$.

	Pertama-tama kita buktikan bahwa funktor ini pada hakikatnya surjektif.
	Tinjau $L \in \Obj(\mathcal{C})$ dan keluarga morfisme
	$(g_j: \alpha H(j) \to L)_{j \in \Obj(J)}$ yang memenuhi syarat
	kompatibilitas $g_{j'} \circ \alpha H(j \to j') = g_j$. Untuk setiap
	$i \in \Obj(I)$, karena $(i/H)$ tidak kosong, kita dapat memilih $j$ dan
	$i \to H(j)$. Definisikan
	$f_i := g_j \circ \alpha(i \to H(j))$. Definisi ini tidak bergantung pada
	pilihan objek koma $(j,i \to H(j))$: memang, jika terdapat
	morfisme-morfisme di $(i/H)$ seperti dalam diagram
	% segment-id: o014.li2.chapter1.filtered-inductive-limits.g003
	\[\begin{tikzcd}[column sep=large]
		& i \arrow[ld] \arrow[d] \arrow[rd] & \\
		H(j) & H(j'') \arrow[l, "{H(j'' \to j)}"] \arrow[r, "{H(j'' \to j')}"'] & H(j')
	\end{tikzcd}\]
	maka dari konstruksinya langsung terlihat bahwa
	% segment-id: o014.li2.chapter1.filtered-inductive-limits.q002
	\[ g_j \circ \alpha(i\to H(j))
	= g_{j''} \circ \alpha(i\to H(j''))
	= g_{j'} \circ \alpha(i\to H(j')). \]
	Karena $(i/H)$ terhubung, ini cukup untuk menunjukkan bahwa $f_i$
	terdefinisi dengan baik. Pembahasan di atas juga menyiratkan bahwa
	$(f_i)_i$ memenuhi syarat kompatibilitas
	$f_i=f_{i'}\circ\alpha(i\to i')$: untuk setiap $i \to i'$ dan
	$i' \to H(j')$, kita dapat menghitung $f_i$ menggunakan objek koma
	$(j',i \to i' \to H(j'))$. Jadi,
	$(L, (f_i)_i) \in \Obj((\alpha/\Delta))$ dipetakan ke
	$(L, (g_j)_j)$.

	Selanjutnya kita tunjukkan bahwa funktor ini penuh dan setia. Memberikan
	suatu morfisme di $(\alpha/\Delta)$, yaitu
	$(L, (f_i)_i) \to (L', (f'_i)_i)$, setara dengan memberikan morfisme di
	$\mathcal{C}$, yaitu $\theta: L \to L'$, sedemikian sehingga
	$\theta f_i = f'_i$ selalu berlaku. Cukup memeriksa syarat ini pada
	objek-objek $H(j)$ dalam citra kofinal $H$: untuk setiap $i$, pilih
	$j \in \Obj(J)$ dan $i \to H(j)$; syarat tersebut kemudian menjadi
	$\theta f_{H(j)} = f'_{H(j)}$, yakni bahwa $\theta$ merupakan morfisme di
	$(\alpha H / \Delta)$.
\end{bukti}

% segment-id: o014.li2.chapter1.filtered-inductive-limits.d004
\begin{definisi}
	\index{kategori!terfilter (filtered)}
	\index{limit!terfilter (filtered)}
	Menurut \cite[Definisi 2.7.6]{Li1}, kategori $I$ dengan
	$\Obj(I)\neq\emptyset$ yang memenuhi syarat-syarat berikut disebut
	\emph{terfilter}\footnote{Catatan penerjemah: sumber tidak menyatakan syarat
	$\Obj(I)\neq\emptyset$, sehingga kategori kosong akan memenuhi kedua butir
	secara vakum; sumber juga menyebut $k\in\Obj(I)$ sebagai morfisme pada butir
	kedua. Target menambahkan syarat ketak-kosongan dan memperbaiki jenis $k$
	menjadi objek (O014-C010 dan O014-C011).}.
	% segment-id: o014.li2.chapter1.filtered-inductive-limits.l003
	\begin{compactitem}
		% segment-id: o014.li2.chapter1.filtered-inductive-limits.i006
		\item Untuk setiap $i, j \in \Obj(I)$, terdapat objek $k \in \Obj(I)$
		beserta morfisme $i \rightarrow k$ dan $j \rightarrow k$.
		% segment-id: o014.li2.chapter1.filtered-inductive-limits.i007
		\item Untuk setiap pasangan morfisme $f, g: i \to j$ di $I$, terdapat
		objek $k \in \Obj(I)$ dan morfisme $h: j \to k$ sedemikian sehingga
		$hf = hg$.
	\end{compactitem}

	Jika $I$ merupakan kategori terfilter dan $\alpha: I \to \mathcal{C}$
	merupakan funktor, maka $\varinjlim \alpha$ yang bersesuaian disebut
	\emph{limit induktif terfilter}.
\end{definisi}

% segment-id: o014.li2.chapter1.filtered-inductive-limits.e002
\begin{contoh}\label{eg:filtered-poset}
	\index{himpunan terurut parsial!terfilter (filtered)}
	Salah satu kelas kategori terfilter yang umum berasal dari himpunan
	terurut parsial terfilter\footnote{Ketika digunakan sebagai indeks limit,
	kategori terfilter sering dapat digantikan dengan himpunan terurut parsial
	terfilter; lihat Catatan \sourcecrossref{rem:filtered-poset}{A.2.3} dan
	penjelasan untuk latihan-latihan dalam
	Lampiran~\sourcecrossref{sec:app-Abel}{A}.}. Setiap himpunan terurut parsial
	$(P, \leq)$ menentukan kategori $\mathcal{P}$ yang bersesuaian sedemikian
	sehingga $x \leq y  \iff \Hom_{\mathcal{P}}(x, y) \neq \emptyset$. Jika
	$\mathcal{P}$ merupakan kategori terfilter, maka $(P, \leq)$ disebut
	\emph{himpunan terurut parsial terfilter}. Suatu himpunan terurut parsial
	bersifat terfilter jika dan hanya jika himpunan itu tidak kosong dan setiap
	dua unsurnya $i, j$ mempunyai batas atas bersama $k$. Setiap himpunan
	terurut total yang tidak kosong jelas terfilter.
\end{contoh}

% segment-id: o014.li2.chapter1.filtered-inductive-limits.pr002
\begin{proposisi}\label{prop:cofinal-filtered}
	Misalkan $I$ kategori terfilter. Maka subkategori penuh $J\subset I$
	bersifat kofinal jika dan hanya jika untuk setiap $i \in \Obj(I)$ terdapat
	$j \in \Obj(J)$ dan morfisme $i \to j$; dalam hal ini $J$ juga terfilter.
\end{proposisi}
% segment-id: o014.li2.chapter1.filtered-inductive-limits.pf002
\begin{bukti}
	Arah ``hanya jika'' jelas. Untuk arah ``jika'', diberikan
	$j \leftarrow i \to j'$, dengan $j, j' \in \Obj(J)$, sifat terfilter
	menjamin adanya diagram komutatif
	% segment-id: o014.li2.chapter1.filtered-inductive-limits.g004
	\[\begin{tikzcd}[row sep=tiny, column sep=small]
		& j \arrow[r] & k & j' \arrow[l] & \\
		j \arrow[equal, ru] & & & & j' \arrow[equal, lu] \\
		& & i \arrow[llu] \arrow[uu] \arrow[rru] & &
	\end{tikzcd}\]
	dengan $k \in \Obj(J)$. Ini menunjukkan bahwa $J$ kofinal.

	Untuk membuktikan bahwa $J$ terfilter, ambil sembarang
	$i, j \in \Obj(J)$. Karena $I$
	terfilter, terdapat $k_0 \in \Obj(I)$ beserta morfisme $i\to k_0$ dan
	$j\to k_0$; pilih pula $k \in \Obj(J)$ beserta
	morfisme $k_0 \to k$. Dengan mengomposisikan morfisme-morfisme tersebut,
	di $J$ terdapat morfisme $i \rightarrow k \leftarrow j$.
	\footnote{Catatan penerjemah: sumber menulis
	$i\rightarrow k_0\leftarrow j$ sebagai morfisme-morfisme di $J$, padahal
	$k_0$ hanya diketahui berada di $I$. Target menggunakan komposit menuju
	$k\in\Obj(J)$ (O014-C012).} Ini membuktikan syarat pertama
	untuk kategori terfilter. Syarat kedua dapat dibuktikan dengan gagasan yang
	sama.
\end{bukti}

% segment-id: o014.li2.chapter1.filtered-inductive-limits.e003
\begin{contoh}
	Definisikan
	$\mathrm{OFin}_I := \left\{ \text{subkategori penuh}\; J \subset I: \Obj(J) \;\text{berhingga} \right\}$.
	Himpunan ini, terhadap $\subset$, merupakan himpunan terurut parsial
	terfilter: untuk setiap $J, K \in \mathrm{OFin}_I$, subkategori yang
	diperoleh dengan mengambil gabungan himpunan objeknya merupakan batas atas
	bersama bagi $J$ dan $K$.
\end{contoh}

% segment-id: o014.li2.chapter1.filtered-inductive-limits.p005
Untuk melihat kegunaan konstruksi ini, tinjau sembarang funktor
$\alpha: I \to \mathcal{C}$ (atau $\beta: I^{\opp} \to \mathcal{C}$).
Pembatasannya pada subkategori penuh $J$ memberikan $\alpha|_J$ (atau
$\beta|_{J^{\opp}}$). Jika $J \subset K$, terdapat morfisme alami
$\varinjlim \alpha|_J \to \varinjlim \alpha|_K$ (atau
$\varprojlim \beta|_{K^{\opp}} \to \varprojlim \beta|_{J^{\opp}}$), asalkan
limit-limit tersebut ada. Hasil berikut menjelaskan cara
menghampiri sembarang limit secara ``terfilter'' melalui kasus dengan
objek berhingga.

% segment-id: o014.li2.chapter1.filtered-inductive-limits.pr003
\begin{proposisi}\label{prop:limit-filtered-approx}
	Dalam situasi di atas, terdapat isomorfisme kanonik
	% segment-id: o014.li2.chapter1.filtered-inductive-limits.q003
	\[ \varinjlim_{\substack{J \in \mathrm{OFin}_I \\ \subset}} \varinjlim \alpha|_J \rightiso \varinjlim \alpha \]
	(atau $\varprojlim \beta \rightiso
	\varprojlim_{J \in \mathrm{OFin}_I^{\opp}}
	\varprojlim \beta|_{J^{\opp}}$), dengan syarat limit rangkap pada ruas kiri
	(atau ruas kanan) ada.\footnote{Catatan penerjemah: pada rumus dual, sumber
	menulis limit induktif luar. Namun untuk $J\subset K$, morfisme transisinya
	berarah dari suku $K$ ke suku $J$, sehingga suku-suku itu membentuk diagram
	pada $\mathrm{OFin}_I^{\opp}$ dan konstruksi luarnya adalah limit projektif
	(O014-C013).}
\end{proposisi}
% segment-id: o014.li2.chapter1.filtered-inductive-limits.pf003
\begin{bukti}
	Ambil kasus $\varinjlim$ sebagai contoh. Sifat universal memberikan bijeksi
	kanonik
	% segment-id: o014.li2.chapter1.filtered-inductive-limits.q004
	\[ \begin{aligned}
	\Hom\left( \varinjlim_{J\in\mathrm{OFin}_I} \varinjlim \alpha|_J, S \right)
	&\simeq \varprojlim_{J\in\mathrm{OFin}_I^{\opp}}
	\Hom\left( \varinjlim \alpha|_J, S \right) \\
	&\simeq \varprojlim_{J\in\mathrm{OFin}_I^{\opp}}
	\varprojlim \Hom(\alpha|_J, S) .
	\end{aligned} \]
	dengan $S \in \Obj(\mathcal{C})$ sembarang, sedangkan ruas kanan merupakan
	$\varprojlim$ rangkap dari himpunan. Cukup diperiksa, untuk keluarga
	himpunan $(X_i)_{i \in \Obj(I)}$ beserta keluarga
	morfisme kompatibel
	$\left( f_{i \to j}: X_j \to X_i \right)_{[i \to j] \in \Mor(I)}$,
	bijeksi
	% segment-id: o014.li2.chapter1.filtered-inductive-limits.g005
	\[\begin{tikzcd}[row sep=tiny]
		\displaystyle\varprojlim_{i \in \Obj(I)} X_i \arrow[r, "\sim"] & \displaystyle\varprojlim_{J \in \mathrm{OFin}_I^{\opp}} \displaystyle\varprojlim_{j \in \Obj(J)} X_j \\
		(x_i)_i \arrow[mapsto, r] & \left( (x_j)_{j \in \Obj(J)} \right)_{J \in \mathrm{OFin}_I} .
	\end{tikzcd}\]

	Untuk setiap $J$, konstruksi konkret $\varprojlim$ dari himpunan ialah
	% segment-id: o014.li2.chapter1.filtered-inductive-limits.q005
	\[ \varprojlim_{j \in \Obj(J)} X_j = \left\{ (x_j)_j \in \prod_{j \in \Obj(J)} X_j : \forall [i \to j] \in \Mor(J), \; f_{i \to j}(x_j) = x_i \right\}, \]
	sehingga bijeksi tersebut jelas: setiap data dalam $\varprojlim_i X_i$,
	baik koordinat $x_i$ maupun syarat kompatibilitas
	$f_{i \to j}(x_j) = x_i$, dapat tercakup oleh suatu
	$J \in \mathrm{OFin}_I$.
\end{bukti}

% segment-id: o014.li2.chapter1.filtered-inductive-limits.p006
Ingat bahwa kategori $\cate{Set}$ mempunyai semua $\varinjlim$ dan
$\varprojlim$ kecil. Limit induktif kecil terfilter mempunyai deskripsi yang
sangat baik.

% segment-id: o014.li2.chapter1.filtered-inductive-limits.pr004
\begin{proposisi}\label{prop:filtered-union}
	Misalkan $I$ kategori kecil yang terfilter dan
	$\alpha: I \to \cate{Set}$ funktor. Definisikan relasi biner berikut pada
	himpunan $\bigsqcup_{i \in \Obj(I)} \alpha(i)$: untuk $x \in \alpha(i)$
	dan $y \in \alpha(j)$, relasi $x \sim y$ berarti bahwa terdapat objek
	$k\in\Obj(I)$ beserta morfisme $i\rightarrow k$ dan $j\rightarrow k$
	sedemikian sehingga
	$\alpha(i \to k)(x) = \alpha(j \to k)(y)$. Maka $\sim$ merupakan relasi
	ekuivalensi, dan
	% segment-id: o014.li2.chapter1.filtered-inductive-limits.q006
	\[ \varinjlim \alpha = \left( \bigsqcup_{i \in \Obj(I)} \alpha(i) \right) \big/ \sim . \]
\end{proposisi}
% segment-id: o014.li2.chapter1.filtered-inductive-limits.pf004
\begin{bukti}
	Lihat pembahasan setelah \cite[Definisi 2.7.6]{Li1}.
\end{bukti}

% segment-id: o014.li2.chapter1.filtered-inductive-limits.p007
Selanjutnya, unsur $\varinjlim \alpha$ di atas akan kita tuliskan sebagai
$[x_{i_0}]$, dengan $i_0 \in \Obj(I)$ dan $[x_{i_0}]$ merupakan kelas
ekuivalensi yang memuat $x_{i_0} \in \alpha(i_0)$. Di sisi lain, untuk
kategori kecil umum $J$ dan funktor $\beta: J^{\opp} \to \cate{Set}$, kita
akan menuliskan unsur $\varprojlim \beta$ dalam bentuk
$(y_j)_{j \in \Obj(J)}$, dengan $y_j \in \beta(j)$ memenuhi syarat
kompatibilitas $\beta(j' \to j)(y_j) = y_{j'}$.

% segment-id: o014.li2.chapter1.filtered-inductive-limits.p008
Berikutnya, tinjau funktor berbentuk
$\alpha: I \times J^{\opp} \to \cate{Set}$, dengan $I, J$ kategori kecil;
menurut peubah $I$ dan $J$ masing-masing dapat diambil $\varinjlim$ dan
$\varprojlim$. Untuk setiap $(i_0, j_0) \in \Obj(I \times J)$, terdapat
komposisi pemetaan alami
% segment-id: o014.li2.chapter1.filtered-inductive-limits.q007
\[ \varprojlim_j \alpha(i_0, j) \to \alpha(i_0, j_0) \to \varinjlim_i \alpha(i, j_0); \]
komposisi ini alami dalam $i_0$ dan $j_0$. Dengan memvariasikan $j_0$,
kita memperoleh pemetaan kanonik
$\varprojlim_j \alpha(i_0, j) \to \varprojlim_j \varinjlim_i \alpha(i, j)$;
selanjutnya, dengan memvariasikan $i_0$, kita memperoleh pemetaan kanonik
% segment-id: o014.li2.chapter1.filtered-inductive-limits.q008
\begin{equation}\label{eqn:filtered-finite-exchange}
	e: \varinjlim_i \varprojlim_j \alpha(i, j) \to \varprojlim_j \varinjlim_i \alpha(i, j).
\end{equation}
Dengan menggunakan notasi wakil, $e$ memetakan
$\left[ (x_{i_0, j})_j \right]$ ke
$\left( [ x_{i_0, j}] \right)_j$.

% segment-id: o014.li2.chapter1.filtered-inductive-limits.p009
Jika $\Mor(J)$ merupakan himpunan berhingga, kategori $J$ disebut berhingga;
hal ini setara dengan mensyaratkan bahwa himpunan objeknya dan himpunan
morfisme di antara setiap dua objek semuanya berhingga. Hasil berikut akan
digunakan dalam \S\sourcecrossref{sec:cat-localization}{1.9}.

% segment-id: o014.li2.chapter1.filtered-inductive-limits.pr005
\begin{proposisi}\label{prop:filtered-finite-exchange}
	Misalkan $I$ kategori kecil yang terfilter, $J$ kategori berhingga, dan
	$\alpha: I \times J^{\opp} \to \cate{Set}$ funktor. Maka pemetaan $e$ dalam
	\eqref{eqn:filtered-finite-exchange} merupakan bijeksi.
\end{proposisi}
% segment-id: o014.li2.chapter1.filtered-inductive-limits.pf005
\begin{bukti}
	Pertama-tama kita tunjukkan bahwa $e$ surjektif. Diberikan
	$(b_j)_{j \in \Obj(J)} \in \varprojlim_j \varinjlim_i \alpha(i, j)$,
	tuliskan setiap $b_j$ sebagai $[x_{i_j, j}]$. Karena $I$ terfilter dan $J$
	berhingga, kita dapat memilih objek sasaran bersama $i\in\Obj(I)$ untuk
	semua $i_j$, lalu mengganti setiap wakil dengan citranya sepanjang morfisme
	$i_j\to i$.
	Dengan demikian, semua wakil dapat ditulis $x_{i,j}$ dengan indeks $i$ yang
	sama. Untuk setiap morfisme $j \to j'$ di $J$, kita mempunyai
	% segment-id: o014.li2.chapter1.filtered-inductive-limits.q009
	\[ \left[ \alpha(i, j \to j') (x_{i, j'}) \right] = \left[ x_{i, j} \right]. \]
	Karena $\Mor(J)$ berhingga dan $I$ terfilter, kita dapat memilih satu
	morfisme $i\to i_1$ yang menyamakan semua pasangan tersebut, mengganti
	semua wakil dengan citranya di indeks $i_1$,
	lalu menamai
	kembali $i_1$ sebagai $i$, sehingga
	$\alpha(i, j \to j') (x_{i, j'}) = x_{i, j}$ berlaku untuk semua
	$[j \to j'] \in \Mor(J)$. Karena itu,
	$e\left( [(x_{i, j})_j] \right) = (b_j)_j$.

	Selanjutnya kita tunjukkan bahwa $e$ injektif. Misalkan $e(x) = e(y)$.
	Karena $I$ terfilter, kita dapat mengganti wakil-wakil $x$ dan $y$ dengan
	citranya di satu indeks bersama $i\in\Obj(I)$, sehingga keduanya mempunyai wakil
	$(x_{i,j})_j$ dan $(y_{i,j})_j$. Jadi, untuk setiap $j \in \Obj(J)$,
	berlaku $[x_{i,j}]=[y_{i,j}]$. Karena $I$ terfilter dan $J$ berhingga, kita
	dapat memilih satu indeks lanjutan bersama, mengganti semua komponen dengan
	citranya di sana, lalu menamainya kembali sebagai $i$, sehingga
	$x_{i,j}=y_{i,j}$ untuk setiap
	$j$. Dengan demikian, $x=y$.
\end{bukti}
